% ==============================================================================
% ABSTRACT
% ==============================================================================
Do next-generation \LLM agents inherit the cooperative biases documented in
their predecessors, or does scale and provider diversity reshape equilibrium
behaviour in competitive multi-agent settings? Willis et al.\ \cite{Willis2025_llm_ipd}
established a benchmark for this question using evolutionary game theory and
the Iterated Prisoner's Dilemma (\IPD), finding consistent cooperative biases
in ChatGPT-4o and Claude~3.5~Sonnet. We extend this benchmark to four frontier
models released in 2025--2026---\textbf{Claude~Sonnet~4.6}, \textbf{Gemini~2.5~Flash},
\textbf{Gemini~3.1~Pro}, and \textbf{GPT-5.4~Mini}---applying the identical
protocol across three prompting styles (Default, Prose, Self-Refine) and four
population compositions (balanced and biased, with and without noise).
Cooperative bias persists across providers (\textbf{H1}): ten of twelve
model--prompt combinations favour cooperative equilibria in balanced noiseless
conditions. Cross-provider divergence is substantial (\textbf{H3}): Gemini~2.5~Flash
reaches up to 77\% aggressive equilibria under biased conditions, while GPT-5.4~Mini
reaches 70\% cooperative equilibria under Self-Refine. Support for aggressive
capability parity is partial (\textbf{H2}): Self-Refine raises \ICD in all
models and Gemini~3.1~Pro Refine achieves the highest \ICD in the dataset
(0.925), but Default and Prose prompts show no systematic narrowing. Noise
robustness improves directionally but not robustly across the Anthropic lineage
(\textbf{H4}): average $\Delta_{\text{noise}}$ falls from 13~pp for
Claude~3.5~Sonnet to 6~pp for Claude~Sonnet~4.6, a gap that is suggestive but
not statistically robust given the predecessor's smaller sample. Provider choice,
not model generation, is the dominant driver of equilibrium outcomes; noise
remains a universal challenge regardless of model size or vintage.
